\documentclass[11pt,a4paper]{article}

\usepackage[T1]{fontenc}
\usepackage[utf8]{inputenc}
\usepackage{lmodern}
\usepackage{microtype}
\usepackage[margin=2.3cm,top=2.6cm,bottom=2.4cm]{geometry}
\usepackage{booktabs}
\usepackage{tabularx}
\usepackage{array}
\usepackage{longtable}
\usepackage{pdflscape}
\usepackage[table]{xcolor}
\usepackage{titlesec}
\usepackage{enumitem}
\usepackage{fancyhdr}
\usepackage{caption}
\usepackage{amssymb}
\usepackage{graphicx}
\usepackage[colorlinks=true,linkcolor=black,urlcolor=black,citecolor=black]{hyperref}

% Monochrome: hierarchy comes from weight, size and rules -- not colour.
\definecolor{accent}{HTML}{000000}
\definecolor{rule}{HTML}{000000}
\definecolor{warn}{HTML}{000000}
\definecolor{ink}{HTML}{000000}
\definecolor{muted}{HTML}{000000}
\color{ink}

\titleformat{\section}
  {\normalfont\Large\bfseries}{\thesection}{0.7em}{}[\vspace{-0.55em}{\color{black}\rule{\linewidth}{0.7pt}}]
\titlespacing*{\section}{0pt}{1.7em}{0.7em}
\titleformat{\subsection}{\normalfont\large\bfseries}{\thesubsection}{0.6em}{}
\titlespacing*{\subsection}{0pt}{1.3em}{0.5em}
\titleformat{\subsubsection}{\normalfont\normalsize\bfseries}{\thesubsubsection}{0.6em}{}

\pagestyle{fancy}
\setlength{\headheight}{22pt}
\addtolength{\topmargin}{-10pt}
\fancyhf{}
\renewcommand{\headrulewidth}{0.3pt}
\fancyhead[L]{\small Submission 2 — Pretrained Models}
\fancyhead[R]{\small Cricket Shot Attention CNN}
\fancyfoot[C]{\small\thepage}

\newenvironment{keybox}
  {\par\medskip\noindent\begin{tabular}{@{}p{0.3em}@{\hspace{0.8em}}p{\dimexpr\linewidth-2.2em}@{}}
   {\color{black}\vrule width 2.5pt} &}
  {\end{tabular}\par\medskip}

\setlist[itemize]{leftmargin=1.3em,itemsep=0.25em,topsep=0.4em}
\setlist[enumerate]{leftmargin=1.5em,itemsep=0.3em,topsep=0.4em}
\renewcommand{\arraystretch}{1.25}
\captionsetup{font=small,labelfont=bf,skip=6pt}
\setlength{\parskip}{0.55em}
\setlength{\parindent}{0pt}

\graphicspath{{figures/}}

\begin{document}

\begin{center}
{\small CNN PROJECT \textbullet\ SUBMISSION 2 OF 4}\\[0.9em]
{\LARGE\bfseries Pretrained Models, With and Without Fine-Tuning}\\[1.0em]
{\large Three ImageNet backbones, two transfer-learning modes, six runs}\\[1.2em]
{\small Repository: \href{https://github.com/Rayhan-who/cricket-shot-attention-cnn}{github.com/Rayhan-who/cricket-shot-attention-cnn}}
\end{center}

\vspace{0.4em}
{\hrule height 0.8pt}
\vspace{1.0em}

\section*{Summary}

Six training runs were conducted: ResNet50, EfficientNetB0 and DenseNet121, each
evaluated as a frozen feature extractor and as a fine-tuned network, for up to 40 epochs
with early stopping. All six completed in 80.5 minutes on a Kaggle T4 GPU.

The best configuration, EfficientNetB0 fine-tuned, reached \textbf{0.9972} test accuracy
and \textbf{0.9973} macro F1 --- 719 of 721 test images correct. Fine-tuning improved
macro F1 for all three architectures, by a mean of $+0.0478$. No run showed
overfitting.

\begin{keybox}
\textbf{These absolute accuracies are almost certainly inflated by residual duplicate
leakage in the dataset, and should not be reported without the qualification set out in
Section~\ref{sec:caveat}.} The relative findings --- the fine-tuning effect, the
per-class ordering and the cost measurements --- are unaffected.
\end{keybox}

\section{Experimental setup}

\subsection{Inherited from Submission 1}

The following were fixed in Submission 1 and are not re-decided here. Holding them
constant is what makes the six runs comparable with each other and with Submissions 3
and 4.

\begin{tabularx}{\linewidth}{@{}>{\bfseries}p{4.6cm}X@{}}
\toprule
Setting & Value \\
\midrule
Split & 70 / 15 / 15, grouped by duplicate identity, stratified by class \\
Split sizes & train 3{,}321 \textbullet\ validation 682 \textbullet\ test 721 \\
Random seed & 42 \\
Channel conversion & forced RGB (4{,}586 of 4{,}724 images are RGBA) \\
Input resolution & $224 \times 224$ \\
Augmentation & rotation $\pm10^{\circ}$, zoom 0.8--1.0, brightness $\pm20\%$, contrast $\pm20\%$ \\
Horizontal flip & \textbf{not applied} --- mirroring inverts the label on leg-side strokes \\
Class weights & not applied --- class imbalance is 1.12\,:\,1 \\
\bottomrule
\end{tabularx}

The split was not transferred as a file. It was \emph{recomputed} from the same
deterministic procedure and then verified against Submission 1's counts; the notebook
asserts the match and reports ``Split reproduces Submission 1 exactly''. An assertion
also confirms that no duplicate group spans a partition boundary.

\subsection{What varies}

\begin{tabularx}{\linewidth}{@{}>{\bfseries}p{4.6cm}X@{}}
\toprule
Factor & Levels \\
\midrule
Backbone & ResNet50, EfficientNetB0, DenseNet121 (ImageNet weights) \\
Transfer mode & frozen (head only) \textbullet\ fine-tuned (whole backbone) \\
Learning rate & $10^{-3}$ frozen, $10^{-4}$ fine-tuned \\
Epoch budget & 40, with early stopping on validation loss, patience 8 \\
\bottomrule
\end{tabularx}

\subsubsection*{Two implementation decisions worth recording}

\textbf{Each backbone's own \texttt{preprocess\_input} is applied inside the model.}
ResNet50 expects Caffe-style mean subtraction, DenseNet Torch-style normalisation, and
EfficientNet raw $[0,255]$. Placing the preprocessing inside each model lets all six runs
share one set of \texttt{uint8} arrays with no possibility of applying the wrong
normalisation --- an error that would be silent and would look like a modelling result.

\textbf{BatchNorm is held in inference mode throughout}, including when fine-tuning. The
backbone is always called with \texttt{training=False}. Updating BatchNorm statistics
while fine-tuning on a small dataset destabilises training and is a common cause of
fine-tuned models scoring \emph{below} frozen ones.

\section{Results}

\subsection{All six runs}

\begin{center}
\small
\begin{tabular}{@{}llrrrrrrr@{}}
\toprule
\textbf{Backbone} & \textbf{Mode} & \textbf{Acc.} & \textbf{Prec.} & \textbf{Rec.} &
\textbf{Macro F1} & \textbf{Ep.} & \textbf{Train} & \textbf{ms/img} \\
\midrule
EfficientNetB0 & fine-tuned & \textbf{0.9972} & 0.9973 & 0.9973 & \textbf{0.9973} & 24 & 10.6\,m & 8.73 \\
DenseNet121    & fine-tuned & 0.9972 & 0.9972 & 0.9973 & 0.9972 & 24 & 17.5\,m & 14.87 \\
ResNet50       & fine-tuned & 0.9945 & 0.9945 & 0.9945 & 0.9944 & 29 & 19.5\,m & 9.75 \\
ResNet50       & frozen     & 0.9750 & 0.9754 & 0.9752 & 0.9752 & 40 & 12.0\,m & 9.72 \\
EfficientNetB0 & frozen     & 0.9598 & 0.9610 & 0.9599 & 0.9602 & 40 & \textbf{6.6\,m} & 9.18 \\
DenseNet121    & frozen     & 0.9098 & 0.9122 & 0.9105 & 0.9103 & 40 & 13.3\,m & 16.47 \\
\bottomrule
\end{tabular}
\normalsize
\end{center}

Precision and recall are macro-averaged. ``Ep.'' is the epoch at which training stopped
out of a budget of 40.

\subsection{The effect of fine-tuning}

\begin{figure}[htbp]
\centering
\includegraphics[width=0.78\linewidth]{03_frozen_vs_finetuned.png}
\caption{Macro F1 for each backbone, frozen versus fine-tuned. The direction is the same
in all three cases; the magnitude is largest where the frozen baseline was weakest.}
\end{figure}

\begin{center}
\begin{tabular}{@{}lrrr@{}}
\toprule
\textbf{Backbone} & \textbf{Frozen} & \textbf{Fine-tuned} & \textbf{Gain} \\
\midrule
DenseNet121    & 0.9103 & 0.9972 & \textbf{+0.0870} \\
EfficientNetB0 & 0.9602 & 0.9973 & +0.0371 \\
ResNet50       & 0.9752 & 0.9944 & +0.0192 \\
\midrule
\textbf{Mean}  &        &        & \textbf{+0.0478} \\
\bottomrule
\end{tabular}
\end{center}

The gain is inversely related to the frozen baseline, which is the expected pattern:
fine-tuning contributes most where the frozen ImageNet representation transfers least
well to the target domain.

A second signal supports the same reading. \textbf{All three frozen runs consumed the
full 40-epoch budget without triggering early stopping}, while all three fine-tuned runs
converged and stopped early --- at 24, 24 and 29 epochs. The frozen models were still
improving when the budget expired; the fine-tuned models had finished.

\subsection{Training curves}

\begin{figure}[htbp]
\centering
\includegraphics[width=\linewidth]{01_training_curves.png}
\caption{Accuracy (top row) and loss (bottom row), training versus validation, for each
of the six runs. Note that the validation curve sits \emph{above} the training curve in
the frozen runs: training accuracy is measured on augmented images and validation on
clean ones, so this ordering is the expected healthy state under this augmentation
policy, not a sign of underfitting.}
\end{figure}

\subsection{Confusion matrices}

\begin{figure}[htbp]
\centering
\includegraphics[width=0.96\linewidth]{02_confusion_matrices.png}
\caption{Test-set confusion matrices for all six runs. Cells show counts; shading is
row-normalised.}
\end{figure}

Submission 1 predicted, from visual inspection of the sample grid, that \texttt{drive}
and \texttt{legglance-flick} would be the hardest pair to separate. The automated check
reports \texttt{drive} $\leftrightarrow$ \texttt{pullshot} instead, so the
\textbf{prediction was not confirmed}.

\begin{keybox}
That result should not be over-read. The best model makes \textbf{two errors in total},
so the ``hardest pair'' is being inferred from a single misclassification. At this
accuracy the confusion matrix carries almost no information. The more reliable signal
comes from the weakest run --- DenseNet121 frozen --- where \texttt{drive} has the lowest
per-class F1 by a clear margin.
\end{keybox}

\subsection{Per-class performance}

\begin{figure}[htbp]
\centering
\includegraphics[width=0.94\linewidth]{04_per_class_f1.png}
\caption{Per-class F1 across all six runs.}
\end{figure}

\begin{center}
\begin{tabular}{@{}lrl@{}}
\toprule
\textbf{Class} & \textbf{Mean F1 across runs} & \\
\midrule
\texttt{drive}           & 0.9601 & hardest \\
\texttt{legglance-flick} & 0.9719 & \\
\texttt{pullshot}        & 0.9752 & \\
\texttt{sweep}           & 0.9826 & easiest \\
\bottomrule
\end{tabular}
\end{center}

The extremes are consistent; the middle is not. \texttt{sweep} is the easiest class in
five of six runs (in ResNet50 fine-tuned it trails \texttt{pullshot} by 0.0002), which
is consistent with it being the most visually distinct posture --- it is the only stroke
played from a lowered or kneeling position, so its silhouette differs from the other three
in a way that survives heavy downsampling. \texttt{drive} is the hardest class in five of
six runs (DenseNet121 fine-tuned is the exception), which is consistent with it being an
upright side-on stroke that shares its body geometry with both \texttt{legglance-flick}
and \texttt{pullshot}. The order of those two middle classes swaps between runs, so only
the mean ordering above should be read as a finding.

\subsection{Cost}

\begin{figure}[htbp]
\centering
\includegraphics[width=\linewidth]{05_cost_vs_accuracy.png}
\caption{Test accuracy against training time (left) and against inference time (right).
Circles are frozen runs, squares fine-tuned.}
\end{figure}

Training ranged from 6.6 minutes (EfficientNetB0 frozen) to 19.5 minutes (ResNet50
fine-tuned). Inference ranged from 8.73 to 16.47 ms per image.

\textbf{EfficientNetB0 fine-tuned is the best configuration on every axis that matters.}
It has the highest accuracy, the fastest inference (8.73 ms/image), the second-fastest
training (10.6 min), and the fewest trainable parameters among the fine-tuned runs
(3.97 M, against 23.49 M for ResNet50). DenseNet121 matches its accuracy but costs 66\%
more training time and 70\% more inference time for no measurable benefit.

\section{Overfitting and underfitting}

\begin{center}
\begin{tabular}{@{}llr l@{}}
\toprule
\textbf{Backbone} & \textbf{Mode} & \textbf{Train--val gap} & \textbf{Assessment} \\
\midrule
ResNet50       & fine-tuned & $+0.0070$ & healthy \\
DenseNet121    & fine-tuned & $+0.0029$ & healthy \\
EfficientNetB0 & fine-tuned & $+0.0012$ & healthy \\
ResNet50       & frozen     & $-0.0416$ & healthy \\
EfficientNetB0 & frozen     & $-0.0433$ & healthy \\
DenseNet121    & frozen     & $-0.1224$ & healthy \\
\bottomrule
\end{tabular}
\end{center}

\textbf{Zero of six runs} exceeded a train--validation accuracy gap of 0.15. No
corrective action was required.

The negative gaps on the frozen runs are expected rather than a warning. Training
accuracy is computed on augmented images while validation uses clean ones, so under this
augmentation policy the validation score sitting above the training score is the normal
state. The largest negative value, $-0.1224$ for DenseNet121 frozen, belongs to the
weakest model overall and reflects a model still improving at the epoch budget rather
than one that has failed to fit.

The regularisation in place throughout was: early stopping on validation loss with
patience 8 and best-weight restoration; \texttt{ReduceLROnPlateau} at factor 0.3; dropout
0.3 before the classifier; rotation, zoom, brightness and contrast augmentation; and
BatchNorm held in inference mode.

\section{A caveat on these numbers}
\label{sec:caveat}

\begin{keybox}
\textbf{The absolute accuracies reported above are almost certainly inflated, and the
inflation is a property of the dataset rather than of the models.}
\end{keybox}

The dataset author states on the Kaggle page that the images were collected from the
internet and then augmented to produce the released collection. Submission 1 controlled
for this by grouping images with identical perceptual hashes and splitting on groups
rather than on files. That removed 608 redundant copies (12.9\%) from consideration and
prevented an estimated 17.9\% of the test set from carrying an exact twin in training.

\textbf{That control is incomplete.} A difference hash is sensitive to geometric change:
a rotated, flipped or cropped copy of a photograph produces a completely different hash
and was therefore treated as an independent image. If the author's augmentation included
any geometric transformation --- and the phrase ``augmented'' usually implies it ---
those variants remain distributed across the train and test partitions, and the models
are being scored partly on photographs they have already seen.

Two observations support this reading.

\begin{enumerate}
\item \textbf{DenseNet121 frozen reached 91.0\% with 4{,}100 trainable parameters.} That
  is a linear classifier on fixed ImageNet features. It cannot learn anything
  cricket-specific. For it to reach 91\% on a fine-grained four-class problem, test
  images must lie very close to training images in feature space.
\item \textbf{99.7\% on a task characterised in Submission 1 as fine-grained.} All four
  classes share one batsman, one pitch, comparable kit and comparable backdrops, with the
  label carried by posture and bat angle across a minority of the frame.
\end{enumerate}

Neither observation is proof. The residual redundancy could genuinely be small and the
task genuinely easy under ImageNet features. Resolving the question requires
embedding-based similarity clustering rather than hashing, because rotated copies remain
close in feature space even when their hashes diverge. That audit is the immediate next
piece of work.

\subsection{What this does and does not affect}

\begin{tabularx}{\linewidth}{@{}p{5.4cm}X@{}}
\toprule
\textbf{Unaffected} & \textbf{Not trustworthy} \\
\midrule
The fine-tuning effect --- the direction held in all three architectures &
Absolute accuracy and macro F1 \\
Per-class extremes --- \texttt{drive} hardest and \texttt{sweep} easiest in five of six runs &
The confusion-matrix analysis at the top of the range \\
Training and inference timings &
Use of 0.9973 as a target for Submission 3 \\
\bottomrule
\end{tabularx}

\subsection{A second limitation, unquantifiable with this dataset}

Beyond duplicate leakage there is \textbf{session correlation}. Photographs taken at the
same match share stadium, kit, lighting and crowd, so a model can learn the match rather
than the shot even when no two images are duplicates. The correct remedy is a
session-wise split, but this dataset ships no match metadata, and the filenames were
renumbered into contiguous sequences after augmentation, destroying any provenance that
might have allowed grouping. This is recorded as a known threat to validity rather than
addressed.

\section{Carrying into Submission 3}

\begin{tabularx}{\linewidth}{@{}>{\bfseries}p{5.0cm}X@{}}
\toprule
Item & Value \\
\midrule
Split & loaded from \texttt{results-data/split\_manifest.csv} --- not regenerated \\
Seed & 42 \\
Input, augmentation, metrics & unchanged \\
Best pretrained macro F1 & 0.9973 (EfficientNetB0 fine-tuned) \\
\bottomrule
\end{tabularx}

Submission 3 differs in exactly one respect: the network is trained from random
initialisation rather than ImageNet weights. A from-scratch CNN is expected to score
below the pretrained results --- roughly 3{,}300 training images is not enough to learn
general visual features from scratch --- and that gap is the measured value of transfer
learning.

\begin{keybox}
\textbf{The figure of 0.9973 should be treated as provisional.} If the leakage audit
shows substantial residual redundancy, this baseline and the Submission 3 result will
both move downwards, and the comparison that matters is the one made on a corrected
split. That also determines whether Submission 4 is viable: an attention module cannot
demonstrate a measurable gain against a baseline already above roughly 96\%, because the
improvement required to clear run-to-run variance would exceed the headroom available.
\end{keybox}

\end{document}
